% Generated by report/tools/tables.py from the result files of the repository. Do not edit.
\begin{table}[tbp]
\centering
\footnotesize
\setlength{\tabcolsep}{2.4pt}
\caption[L'année et les lecteurs démêlés (écart 36)]{Écart 36 de \file{study/PROTOCOL.md} (écrit avant toute lecture et tout calcul) : les 98 articles expérimentaux de la comparaison principale (48 de 2018, 50 de 2026) relus par huit agents neufs, T1 à T8, qui lisent chacun les deux années (6 articles de 2018 et 6 ou 7 de 2026). Différence 2026 moins 2018 en points, intervalle de Newcombe à 95~\%, test exact de Fisher avec Holm sur les quatre critères. Effet de l'année : modèle logistique mixte par critère, logit P(oui) = $b_0 + b_1 \times [2026] + u$, $u$ effet aléatoire du lecteur d'écart type $s$, maximum de vraisemblance par quadrature de Gauss-Hermite ; rapport de cotes $e^{b_1}$ avec son intervalle de Wald à 95~\%, test du rapport de vraisemblance de $b_1 = 0$ avec Holm sur les quatre critères ; $s$ avec son test à la frontière (demi-khi-deux). Règle fixée d'avance : un verdict tient si les deux tests passent Holm à 5~\% avec un effet positif ; « même sens, non établi » si les deux effets sont positifs et qu'un test échoue. \textbf{En gras}, les verdicts qui tiennent. Chaque intervalle de Newcombe recouvre celui de la comparaison principale (écarts d'origine : Q5 +43 ; Y9 +36 ; X7 +31 ; X2 +26 points). Dernière ligne : la lecture où toute page de projet compte « oui » pour Q5, la moins favorable à son verdict. Aucune décision invalide. Source : \file{study/r36\_stats.json} (\file{study/r36\_stats.py}).}
\label{tab:r36}
\begin{tabularx}{\linewidth}{@{}l>{\raggedright\arraybackslash}Xrrrrrrrrr>{\raggedright\arraybackslash}p{1.9cm}@{}}
\toprule
 & & \multicolumn{2}{c}{2018} & \multicolumn{2}{c}{2026} & \multicolumn{2}{c}{Différence, Fisher} & \multicolumn{2}{c}{Effet de l'année} & &  \\
\cmidrule(lr){3-4}\cmidrule(lr){5-6}\cmidrule(lr){7-8}\cmidrule(lr){9-10}
Code & Critère & $k/n$ & \% & $k/n$ & \% & \makecell[r]{points\\{}[Newcombe]} & \makecell[r]{$p$\\{}Holm} & \makecell[r]{rapport de cotes\\{}[Wald]} & \makecell[r]{$p$\\{}Holm} & \makecell[r]{$s$ lecteurs\\{}($p$)} & Verdict \\
\midrule
\textbf{Q5} & \textbf{Code en lien} & 15/48 & 31 & 36/50 & 72 & \makecell{+41\\{}[+21,1\,;\,+56,2]} & 0,0004 & \makecell{5,66\\{}[2,37\,;\,13,48]} & 0,0002 & 0,00 (1) & tient \\
\textbf{Y9} & \textbf{Sensibilité à un hyperparamètre} & 20/48 & 42 & 39/50 & 78 & \makecell{+36\\{}[+17,0\,;\,+52,1]} & 0,001 & \makecell{4,96\\{}[2,06\,;\,11,99]} & 0,0006 & 0,00 (1) & tient \\
\textbf{X7} & \textbf{Calcul chiffré} & 15/48 & 31 & 33/50 & 66 & \makecell{+35\\{}[+15,0\,;\,+50,9]} & 0,001 & \makecell{4,37\\{}[1,83\,;\,10,42]} & 0,001 & 0,24 (0,39) & tient \\
X2 & Incertitude rapportée & 33/48 & 69 & 43/50 & 86 & \makecell{+17\\{}[+0,6\,;\,+33,0]} & 0,053 & \makecell{2,88\\{}[1,03\,;\,8,03]} & 0,037 & 0,44 (0,25) & même sens, non établi \\
\quad Q5 & toute page de projet comptée oui & 20/48 & 42 & 37/50 & 74 & \makecell{+32\\{}[+12,8\,;\,+48,7]} & 0,004 &  &  &  &  \\
\bottomrule
\end{tabularx}
\end{table}
